% ch39_realization_nerve_kan.tex — Volume IV Chapter 3
% Retrofitted with full Vol I-III pedagogy: cycle stages, etymology, dialogs.

\chapter{Geometric Realization, the Nerve, and Kan Complexes}

\begin{overviewbox}
Two adjunctions package the relationship between simplicial sets and the
two ``natural'' categories they look like: topological spaces and
ordinary categories.

The first is \term{geometric realization} $\lvert{-}\rvert :
\mathbf{sSet} \to \mathbf{Top}$, left adjoint to the singular complex
functor $\mathrm{Sing}$ of Chapter~38. It is the left Kan extension of
the cosimplicial space $\lvert\Delta^\bullet\rvert$ along the Yoneda
embedding — concretely, the coend $\lvert X\rvert = \int^{[n]} X_n \cdot
\lvert\Delta^n\rvert$. Under this adjunction, simplicial sets serve as a
combinatorial model for topological spaces up to weak equivalence.

The second is the \term{nerve} $\mathrm{N} : \mathbf{Cat} \to \mathbf{sSet}$,
right adjoint to the fundamental category functor $\tau_1 : \mathbf{sSet}
\to \mathbf{Cat}$. The nerve is fully faithful, and a theorem of
Grothendieck identifies its image: the simplicial sets in which every
\emph{inner} horn has a \emph{unique} filler. Relaxing ``unique'' to
``exists'' gives quasi-categories (Chapter~43); insisting on fillers for
\emph{all} horns gives Kan complexes — the topic of the final section.

Kan complexes are the simplicial-set avatars of $\infty$-groupoids:
spaces whose homotopical content is fully captured by simplicial
combinatorics. By the end of the chapter we have the picture
\begin{equation*}
  \text{ordinary categories} \;\hookleftarrow\; \text{simplicial sets} \;\hookrightarrow\;
  \text{Kan complexes} \;\leftrightsquigarrow\; \text{topological spaces},
\end{equation*}
which Chapters~40--42 refine via model-categorical machinery.
\end{overviewbox}


% ═══════════════════════════════════════════════════════════════════════════
\section{Geometric Realization}
\label{sec:geometric-realization}
% ═══════════════════════════════════════════════════════════════════════════

\subsection{Informally}

Every simplicial set $X$ has a topological realization $\lvert X\rvert$
built by gluing one copy of the standard topological $n$-simplex
$\lvert\Delta^n\rvert$ for each non-degenerate $n$-simplex $x \in X_n$,
along the gluing prescribed by $X$'s face operators. Degenerate simplices
contribute nothing new: they are already present as faces of lower
non-degenerate simplices.

In the other direction (Chapter~38), every topological space $T$ has a
singular simplicial set $\mathrm{Sing}(T)$ whose $n$-simplices are
continuous maps $\lvert\Delta^n\rvert \to T$. The two operations are
adjoint — and not by accident: realization is the left Kan extension
along the Yoneda embedding $\Delta \hookrightarrow \mathbf{sSet}$ of the
cosimplicial space of standard simplices.

\begin{readernote}
A useful way to picture realization: imagine the simplicial set as a
``set of instructions'' for assembling a topological space.
\begin{itemize}
  \item ``Take this many vertices.''
  \item ``Glue these vertices into edges according to the face data.''
  \item ``Glue these edges into triangles, these triangles into
        tetrahedra, etc.''
\end{itemize}
The realization is the topological space the instructions describe.
Crucially, only \emph{non-degenerate} simplices add new content;
degeneracies are bookkeeping.
\end{readernote}

\subsection{Expanding}

The coend formula $\lvert X\rvert = \int^{[n]} X_n \cdot \lvert\Delta^n\rvert$
hides a familiar construction. Unrolling:

\begin{itemize}
  \item For each $n$, take a disjoint union of $\lvert X_n\rvert$-many
        copies of $\lvert\Delta^n\rvert$: this is $X_n \cdot
        \lvert\Delta^n\rvert$ (the \emph{copower}).
  \item Glue along face and degeneracy maps: identify a face of an
        $n$-simplex copy with the corresponding $(n-1)$-simplex copy,
        and collapse degenerate copies onto their underlying
        non-degenerate ones.
\end{itemize}
The coend is the universal way to perform all these identifications at once.

\begin{nameorigin}
The \term{coend} (notation $\int^{[n]} \cdots$) is a categorical
construction introduced by Yoneda and developed by Day–Kelly (1969). The
\emph{co-} prefix marks it as the dual of the \emph{end} ($\int_{[n]}$,
a limit-style construction). The integral notation is a metaphor: just
as $\int f$ accumulates values of $f$ across the domain, $\int^{[n]} F$
accumulates the values of $F([n], [n])$ in a categorically-correct way.
The notation has stuck because the analogies with classical integration
(Fubini's theorem, etc.) carry over surprisingly well to the categorical
setting.
\end{nameorigin}

\subsection{Formalizing}

\begin{definition}[Geometric realization]
\label{def:geometric-realization}
The \term{geometric realization} of a simplicial set $X$ is the
topological space
\begin{equation*}
  \lvert X\rvert \;=\; \int^{[n] \in \Delta} X_n \cdot \lvert\Delta^n\rvert,
\end{equation*}
where $S \cdot T = S \times T$ for a set $S$ and space $T$ (the copower
in $\mathbf{Top}$). This defines a functor $\lvert{-}\rvert : \mathbf{sSet}
\to \mathbf{Top}$.
\end{definition}

\begin{howtosay}
$\lvert X\rvert$ \sayit{``the realization of $X$''}; the coend formula
\sayit{``integrate over $\Delta$, copowering $X_n$ by the topological
$n$-simplex''}. The integral notation reflects that the construction
identifies $X_m \cdot \lvert\Delta^m\rvert$ and $X_n \cdot
\lvert\Delta^n\rvert$ along all face/degeneracy maps (Chapter~26 coend
calculus).
\end{howtosay}

\begin{theorem}[Realization–singular adjunction]
\label{thm:realization-sing-adj}
The realization functor is left adjoint to the singular simplicial set
functor:
\begin{equation*}
  \lvert{-}\rvert \;\dashv\; \mathrm{Sing} : \mathbf{Top} \to \mathbf{sSet}, \quad
  \mathbf{Top}(\lvert X\rvert, T) \;\cong\; \mathbf{sSet}(X, \mathrm{Sing}\, T).
\end{equation*}
\end{theorem}

\begin{proof}[Proof sketch]
Both sides are continuous in $X$. By the coend formula and the
$\hom$-coend adjunction of Chapter~26,
\begin{equation*}
  \mathbf{Top}\Bigl(\int^{[n]} X_n \cdot \lvert\Delta^n\rvert, T\Bigr)
  \;\cong\; \int_{[n]} \mathbf{Top}\bigl(X_n \cdot \lvert\Delta^n\rvert, T\bigr)
  \;\cong\; \int_{[n]} \mathbf{Set}\bigl(X_n, \mathbf{Top}(\lvert\Delta^n\rvert, T)\bigr).
\end{equation*}
The rightmost end is $\mathrm{Nat}(X_\bullet, \mathrm{Sing}(T)_\bullet) =
\mathbf{sSet}(X, \mathrm{Sing}\, T)$.
\qed
\end{proof}

\begin{remark}[Realization is a left Kan extension]
\label{rem:realization-lkan}
Equivalently: $\lvert{-}\rvert$ is the left Kan extension of
$\lvert\Delta^\bullet\rvert : \Delta \to \mathbf{Top}$ along the Yoneda
embedding $y : \Delta \to \mathbf{sSet}$,
\begin{equation*}
  \lvert{-}\rvert \;=\; \operatorname{Lan}_y \lvert\Delta^\bullet\rvert.
\end{equation*}
This is the general construction: a cosimplicial object in any cocomplete
category $\mathcal{D}$ extends along Yoneda to a colimit-preserving
functor $\mathbf{sSet} \to \mathcal{D}$ (Chapter~11/18 Kan extension theory).
\end{remark}

\begin{readernote}
The adjunction $\lvert{-}\rvert \dashv \mathrm{Sing}$ is the precise
bridge between combinatorial and topological category theory. The
unit and counit are not isomorphisms in general — there are simplicial
sets that don't recover from their realization, and topological spaces
that don't recover from their singular complex on the nose — but they
\emph{are} weak equivalences after suitable replacement (Chapter~42's
homotopy hypothesis).
\end{readernote}

\subsection{Formally}

\begin{example_thm}[Realization in low dimension]
\label{ex:realization-low}
$\lvert\Delta^n\rvert$ in the abstract is the topological $n$-simplex
already encountered (\S37.\ref{sec:geometric-realization-preview});
$\lvert\partial\Delta^n\rvert$ is the topological boundary sphere $S^{n-1}$
(up to homeomorphism); $\lvert\Lambda^n_k\rvert$ is a contractible piece
of the boundary obtained by removing one $(n-1)$-face.

The simplicial circle $S^1 = \Delta^1/\partial\Delta^1$
(Example~38.\ref{ex:simplicial-circle}) realizes to the topological circle.
\end{example_thm}

\begin{remark}[A common topological subtlety]
The codomain $\mathbf{Top}$ in Theorem~\ref{thm:realization-sing-adj}
should strictly be a category of \emph{nice} spaces — compactly generated
weak Hausdorff is standard — to ensure colimits behave well. In
Chapter~41 we discuss the model-category framework that makes this
precise; for our purposes any cartesian-closed category of spaces with
the standard simplices will do.
\end{remark}


% ═══════════════════════════════════════════════════════════════════════════
\section{The Nerve of a Category}
\label{sec:nerve}
% ═══════════════════════════════════════════════════════════════════════════

\subsection{Informally}

The second adjunction lives between simplicial sets and ordinary
categories. Every ordinary category $\mathcal{C}$ has a simplicial-set
shadow — its nerve $\mathrm{N}(\mathcal{C})$ — whose $n$-simplices are
the length-$n$ composable chains in $\mathcal{C}$:
\begin{equation*}
  c_0 \xrightarrow{f_1} c_1 \xrightarrow{f_2} \cdots \xrightarrow{f_n} c_n.
\end{equation*}
The face operators perform either dropping (at the ends) or composing
(in the middle); the degeneracy operators insert identity arrows.

The construction is fully faithful: $\mathbf{Cat} \hookrightarrow
\mathbf{sSet}$ via the nerve. Grothendieck's theorem identifies the
image \emph{exactly}: a simplicial set is in the image of the nerve iff
every inner horn extends \emph{uniquely}. Relax ``uniquely'' to ``at
least one'' and the resulting class is the quasi-categories.

\subsection{Formalizing}

A finite ordinal $[n] = \{0 < 1 < \cdots < n\}$ is itself an ordinary
category: objects are the elements; there is a unique morphism $i \to j$
whenever $i \leq j$, and composition is forced. This gives a functor
\begin{equation*}
  [-]: \Delta \longrightarrow \mathbf{Cat}, \qquad [n] \longmapsto [n] \text{ as a poset}.
\end{equation*}

\begin{definition}[Nerve of a category]
\label{def:nerve}
For a small category $\mathcal{C}$, the \term{nerve} $\mathrm{N}(\mathcal{C})$
is the simplicial set
\begin{equation*}
  \mathrm{N}(\mathcal{C})_n \;=\; \mathbf{Cat}\bigl([n], \mathcal{C}\bigr).
\end{equation*}
An $n$-simplex is a functor $[n] \to \mathcal{C}$, equivalently an $n$-fold
composable chain $c_0 \xrightarrow{f_1} c_1 \xrightarrow{f_2} \cdots
\xrightarrow{f_n} c_n$ of morphisms in $\mathcal{C}$.
\end{definition}

\begin{howtosay}
$\mathrm{N}(\mathcal{C})$ \sayit{``the nerve of $\mathcal{C}$''}; the
chain description \sayit{``length-$n$ string of composable arrows in
$\mathcal{C}$''}. Some texts write $\mathrm{N}_\bullet(\mathcal{C})$ or
$\mathrm{Ner}(\mathcal{C})$; the bullet ${}_\bullet$ emphasizes
simpliciality but is omitted when context is clear.
\end{howtosay}

\begin{nameorigin}
\term{Nerve} (recall from Chapter~38): biological metaphor introduced by
Grothendieck around 1960 and developed by Graeme Segal (1968). Just as
a nervous system transmits information through an organism, the nerve of
a category transmits its categorical content into the world of
simplicial sets — letting the tools of simplicial homotopy theory
operate on categorical structure.
\end{nameorigin}

\begin{remark}[Face and degeneracy on the nerve]
For a chain $c_0 \to c_1 \to \cdots \to c_n$:
\begin{itemize}
  \item $d_0$ drops $c_0$ (and the morphism $f_1$): gives $c_1 \to \cdots
        \to c_n$.
  \item $d_n$ drops $c_n$ (and $f_n$): gives $c_0 \to \cdots \to c_{n-1}$.
  \item For $0 < i < n$, $d_i$ composes the two adjacent morphisms
        $f_{i+1} \circ f_i$ to give a length-$(n-1)$ chain.
  \item $s_i$ inserts an identity morphism after position $i$, lengthening
        the chain.
\end{itemize}
\end{remark}

\begin{readernote}
The asymmetry between $d_0$ / $d_n$ (which drop) and $d_i$ for $0 < i < n$
(which compose) is the geometric reason that \emph{inner horns} are
special. Filling an inner horn $\Lambda^n_k$ ($0 < k < n$) asks: ``given
all the faces except $d_k$, can we recover $d_k$ and the whole
$n$-simplex?'' For nerves, $d_k$ is forced — it is the composite of two
adjacent arrows. For an outer horn ($k = 0$ or $k = n$), $d_k$ is a
\emph{drop}, which is not forced by the other faces (you would need an
inverse of a morphism, which doesn't generally exist).
\end{readernote}

\begin{theorem}[Fundamental category–nerve adjunction]
\label{thm:nerve-tau1-adj}
The nerve functor $\mathrm{N} : \mathbf{Cat} \to \mathbf{sSet}$ has a
left adjoint $\tau_1 : \mathbf{sSet} \to \mathbf{Cat}$, the
\term{fundamental category} (or \emph{homotopy category}) functor:
\begin{equation*}
  \tau_1 \;\dashv\; \mathrm{N}, \qquad
  \mathbf{Cat}(\tau_1 X, \mathcal{C}) \;\cong\; \mathbf{sSet}(X, \mathrm{N}(\mathcal{C})).
\end{equation*}
$\tau_1 X$ has objects $X_0$ and is freely generated by the $1$-simplices
of $X$ subject to the relations imposed by the $2$-simplices (each
$2$-simplex says ``$f_2 \circ f_1 = f_3$'' for its three boundary
$1$-simplices).
\end{theorem}

\begin{nameorigin}
The notation \term{$\tau_1$} comes from \emph{truncation}: $\tau_1$ takes
a simplicial set (which has data in every dimension) and produces a
$1$-categorical \emph{truncation} of it — discarding the higher
homotopical information and keeping only the categorical part. The
subscript $1$ marks the target dimension (an ordinary $1$-category).
Generalizations $\tau_n$ produce $n$-categorical truncations (Chapter~44).
\end{nameorigin}

\begin{proof}[Proof sketch]
$\tau_1 = \operatorname{Lan}_y [-]$, the left Kan extension along Yoneda
of the cosimplicial category $[-] : \Delta \to \mathbf{Cat}$ — exactly
parallel to the realization construction of
Remark~\ref{rem:realization-lkan}. Adjointness follows from the Yoneda
adjunction for Kan extensions.
\qed
\end{proof}

\subsection{Reorganizing: characterizing the nerve}

\begin{nameorigin}
The \term{characterization of nerves via unique inner-horn filling} is
attributed to \emph{Alexander Grothendieck} (1928--2014), one of the
most influential mathematicians of the 20th century. The result appears
implicitly in Grothendieck's work on classifying spaces and explicitly
in subsequent literature on Segal's nerve construction. Grothendieck's
broader project of reorganizing algebraic geometry through category
theory pervades Volume~IV: the Grothendieck construction (Chapter~46
fibrations), Grothendieck universes, Grothendieck topologies (Chapter~58
DAG), and the homotopy hypothesis (\S\ref{sec:homotopy-hypothesis}
below) are all his.
\end{nameorigin}

\begin{theorem}[Grothendieck's characterization of the nerve]
\label{thm:nerve-characterization}
The nerve functor $\mathrm{N} : \mathbf{Cat} \to \mathbf{sSet}$ is fully
faithful. Its essential image consists of the simplicial sets $X$ in
which every inner horn admits a \emph{unique} filler: for each $0 < k
< n$ and each map $\Lambda^n_k \to X$, there is exactly one $\Delta^n
\to X$ extending it.
\end{theorem}

\begin{proof}[Proof sketch]
\emph{Fully faithful.} A natural transformation $\mathrm{N}(F) \to
\mathrm{N}(G)$ is determined by its action on $0$- and $1$-simplices —
i.e., by an assignment to objects and morphisms compatible with
identities and composition — and that is the same as a functor
$F \to G$.

\emph{Image.} If $X = \mathrm{N}(\mathcal{C})$, an inner horn
$\Lambda^n_k \to X$ for $0 < k < n$ supplies a chain of $n-1$ morphisms
with a ``gap'' at position $k$; the unique filler is forced because
composition in $\mathcal{C}$ is uniquely defined. Conversely, if all
inner horns fill uniquely, the unique-composite $\Lambda^2_1 \to X$
extension defines a composition law on $X_1$ that is associative (forced
by $\Lambda^3_1$ and $\Lambda^3_2$ fillings) and unital (forced by
degeneracies), making $X$ the nerve of a category.
\qed
\end{proof}

\begin{remark}[The quasi-category relaxation]
Theorem~\ref{thm:nerve-characterization} is the gateway to
$\infty$-category theory. \emph{Drop ``unique''} from the
characterization — require inner horns to admit at least one filler, not
necessarily unique — and the resulting class of simplicial sets is the
\term{quasi-categories} (Chapter~43). Composition becomes existential
(a composite exists) rather than constructive (a composite is given);
associativity holds up to higher coherent homotopies rather than on the
nose; and the entire structure of ordinary category theory generalizes
to a ``coherent'' version.
\end{remark}


% ═══════════════════════════════════════════════════════════════════════════
\section{Kan Complexes}
\label{sec:kan-complexes}
% ═══════════════════════════════════════════════════════════════════════════

\subsection{Informally}

A Kan complex is a simplicial set in which \emph{every} horn — inner or
outer — admits a filler. The condition is symmetric: morphisms can be
composed both forward (inner-horn filling, as for the nerve) and
``inverted'' or ``transported backward'' (outer-horn filling).

In an ordinary category nerve, outer horns do \emph{not} fill in general:
$\Lambda^2_0 \to \mathrm{N}(\mathcal{C})$ is two arrows out of a common
source, and a $\Delta^2$ filler would force the existence of an arrow
relating the two targets — which exists only if the category has
limits/colimits or all morphisms are isomorphisms. Outer horns force
invertibility; that is why Kan complexes are sometimes called
\term{$\infty$-groupoids}.

\begin{nameorigin}
\term{Kan complex} is named for \emph{Daniel M.\ Kan} (1927--2013),
Israeli-American mathematician primarily at MIT. Kan introduced the
horn-filling condition in his 1958 paper ``On c.s.s. complexes'' (\emph{c.s.s.}
= ``complete semi-simplicial,'' an older term that has since been
superseded by ``simplicial set''). Kan's other major contributions
include the theory of Kan extensions (Chapters~11, 18, 31, 51), adjoint
functors (Chapter~7), and the Kan model structure on $\mathbf{sSet}$
(Chapter~41).
\end{nameorigin}

\subsection{Formalizing}

\begin{definition}[Kan complex]
\label{def:kan-complex}
A simplicial set $X$ is a \term{Kan complex} if every horn $\Lambda^n_k
\to X$ (for $n \geq 1$ and $0 \leq k \leq n$) extends to $\Delta^n \to X$.
\end{definition}

\begin{howtosay}
``Kan complex'' \sayit{``Kan complex'' or ``Kan simplicial set''}, after
Daniel Kan. The term ``fibrant simplicial set'' is also standard once
Chapter~41's model structure is in place — Kan complexes are exactly
the fibrant objects of the standard (Kan–Quillen) model structure on
$\mathbf{sSet}$.
\end{howtosay}

\begin{nameorigin}
The term \term{$\infty$-groupoid} dates to Grothendieck's \emph{Pursuing
Stacks} manuscript (1983, unpublished but widely circulated). The word
\term{groupoid} itself was coined by Heinrich Brandt (1927) for a
generalization of groups in which composition is only partially defined.
A groupoid has all morphisms invertible; an $\infty$-groupoid has
``all morphisms invertible at every level'' — including the homotopies
between morphisms, the homotopies between those homotopies, and so on
ad infinitum. Kan complexes are the simplicial-set realization of this
``invertibility-at-every-level'' picture.
\end{nameorigin}

\begin{readernote}
The slogan ``Kan complex = $\infty$-groupoid'' is one of the most
important slogans of higher category theory. It says: the simplicial
combinatorial structure where ``all horns fill'' is exactly the
homotopy-theoretic structure where ``everything is invertible.'' The
outer-horn condition encodes the invertibility; the inner-horn condition
encodes the composition. Removing only the outer-horn requirement gives
quasi-categories — $\infty$-categories where invertibility is not forced.
\end{readernote}

\begin{theorem}[$\mathrm{Sing}(T)$ is a Kan complex]
\label{thm:sing-kan}
For any topological space $T$, the singular simplicial set
$\mathrm{Sing}(T)$ is a Kan complex.
\end{theorem}

\begin{proof}[Proof sketch]
A horn $\Lambda^n_k \to \mathrm{Sing}(T)$ corresponds, by adjunction
(Theorem~\ref{thm:realization-sing-adj}) and the unit on representables,
to a continuous map $\lvert\Lambda^n_k\rvert \to T$. The space
$\lvert\Lambda^n_k\rvert$ is a deformation retract of
$\lvert\Delta^n\rvert$ — collapsing the missing face inward gives the
retraction — and any retract extends maps along itself, so
$\lvert\Lambda^n_k\rvert \to T$ extends to $\lvert\Delta^n\rvert \to T$.
Apply the adjunction in reverse.
\qed
\end{proof}

\subsection{The homotopy hypothesis}
\label{sec:homotopy-hypothesis}

\begin{nameorigin}
The \term{homotopy hypothesis} was formulated by \emph{Alexander
Grothendieck} in his 1983 manuscript \emph{Pursuing Stacks} (\emph{À la
poursuite des champs}). It is the thesis that the homotopy theory of
topological spaces and the homotopy theory of $\infty$-groupoids are
\emph{the same thing} — i.e., $\infty$-groupoids are a complete model
for homotopy types of spaces. The hypothesis ramifies through all of
modern algebraic topology and provides the philosophical foundation for
homotopy type theory (Chapter~59).
\end{nameorigin}

\begin{remark}[The homotopy hypothesis: a preview]
The functors $\lvert{-}\rvert$ and $\mathrm{Sing}$ are not equivalences
of categories: they restrict to a \emph{Quillen equivalence} between the
model category of simplicial sets (Kan model structure, Chapter~41) and
the model category of topological spaces (Quillen model structure).
Localizing both sides at weak equivalences produces equivalent homotopy
categories.

Within $\mathbf{sSet}$, the full subcategory of Kan complexes models the
homotopy category of spaces. This statement — that ``$\infty$-groupoids
$=$ homotopy types'' — is Grothendieck's \term{homotopy hypothesis},
formalized within the Kan-complex framework.
\end{remark}

\subsection{Formally: nerve vs.\ Kan vs.\ quasi-category}

\begin{remark}[The three horn conditions]
\label{rem:three-horn-conditions}
The three horn conditions partition the simplicial-set landscape:
\begin{itemize}
  \item \term{Nerve of a category}: every inner horn fills \emph{uniquely}.
  \item \term{Quasi-category}: every inner horn fills (existence only).
  \item \term{Kan complex}: every horn — inner \emph{or outer} — fills.
\end{itemize}
The nerve condition is sharpest; the quasi-category condition relaxes
uniqueness; the Kan condition adds outer-horn filling. Every nerve of a
groupoid is a Kan complex; every Kan complex is a quasi-category; and
both ordinary categories and topological spaces embed into the
simplicial landscape under these descriptions.
\end{remark}

\begin{readernote}
A useful diagram of the simplicial-set landscape:
\begin{itemize}
  \item \emph{Nerves of categories} (inner horns fill \emph{uniquely}) —
        the rigid case.
  \item \emph{Quasi-categories} (inner horns fill, possibly non-uniquely) —
        the $\infty$-category case (Ch.~43).
  \item \emph{Kan complexes} (all horns fill, possibly non-uniquely) —
        the $\infty$-groupoid case.
  \item Intersection: \emph{nerves of groupoids} — both rigid and
        invertible.
\end{itemize}
The whole of $\infty$-category theory lives in the space between these
conditions, and the rest of Volume~IV explores it.
\end{readernote}

\begin{example_thm}[Nerves and groupoids]
\label{ex:nerve-groupoid}
If $\mathcal{G}$ is a groupoid (every morphism invertible), then
$\mathrm{N}(\mathcal{G})$ is a Kan complex: outer horns fill because the
required morphism exists and is unique. Conversely, a Kan complex in
which every inner horn fills \emph{uniquely} is the nerve of a groupoid.
The fully faithful embedding $\mathbf{Grpd} \hookrightarrow \mathbf{sSet}$
lands in the intersection ``nerve $\cap$ Kan.''
\end{example_thm}


% ═══════════════════════════════════════════════════════════════════════════
\section{Exercises}
\label{sec:realization-nerve-exercises}
% ═══════════════════════════════════════════════════════════════════════════

\begin{exercisesbox}
\begin{qa}
\qaitem{Define the geometric realization $\lvert X\rvert$.}{$\lvert X\rvert = \int^{[n]} X_n \cdot \lvert\Delta^n\rvert$ — the coend that glues $X_n$ copies of the topological $n$-simplex along all face and degeneracy identifications.}
\qaitem{What adjunction does realization fit into?}{$\lvert{-}\rvert \dashv \mathrm{Sing}$: $\mathbf{Top}(\lvert X\rvert, T) \cong \mathbf{sSet}(X, \mathrm{Sing}\, T)$.}
\qaitem{Realization as a Kan extension?}{$\lvert{-}\rvert = \operatorname{Lan}_y \lvert\Delta^\bullet\rvert$, the left Kan extension along $y : \Delta \hookrightarrow \mathbf{sSet}$ of the cosimplicial space of standard simplices.}
\qaitem{Where does the coend notation $\int^{[n]}$ come from?}{Yoneda and Day–Kelly (1969). The \emph{co-} prefix marks it as the dual of the end $\int_{[n]}$; the integral metaphor captures categorical-accumulation analogous to classical integration.}
\qaitem{What is $\lvert\partial\Delta^n\rvert$?}{Homeomorphic to the $(n-1)$-sphere $S^{n-1}$.}
\qaitem{What is $\lvert\Lambda^n_k\rvert$?}{A deformation retract of $\lvert\Delta^n\rvert$ — geometrically, the $n$-simplex with one $(n-1)$-face removed; contractible.}
\qaitem{Define the nerve $\mathrm{N}(\mathcal{C})$.}{$\mathrm{N}(\mathcal{C})_n = \mathbf{Cat}([n], \mathcal{C})$, where $[n]$ is the ordinal $0 \to 1 \to \cdots \to n$ regarded as a category. Equivalently, $n$-fold composable chains in $\mathcal{C}$.}
\qaitem{Why ``nerve''?}{Biological metaphor: as a nervous system transmits information through an organism, the nerve of a category transmits its categorical content into $\mathbf{sSet}$, letting simplicial-homotopy tools operate on it.}
\qaitem{What does $d_i$ do on a chain $c_0 \to \cdots \to c_n$ in the nerve?}{$d_0$ drops the leftmost arrow; $d_n$ drops the rightmost; $d_i$ for $0 < i < n$ composes the $i$-th and $(i+1)$-th arrows.}
\qaitem{What does $s_i$ do?}{Inserts an identity morphism on $c_i$, lengthening the chain by one.}
\qaitem{What is the fundamental category $\tau_1 X$?}{The category with objects $X_0$, freely generated by $1$-simplices subject to relations from $2$-simplices. $\tau_1 \dashv \mathrm{N}$.}
\qaitem{Where does the notation $\tau_1$ come from?}{From \emph{truncation}: $\tau_1$ truncates a simplicial set to its $1$-categorical content (discarding higher coherences). The subscript $1$ marks the target dimension.}
\qaitem{Why is $\mathrm{N}$ fully faithful?}{Natural transformations $\mathrm{N}(F) \to \mathrm{N}(G)$ are determined by their $1$-simplex action, which is exactly a functor $F \to G$.}
\qaitem{State Grothendieck's characterization of the nerve's image.}{$X$ is in the essential image of $\mathrm{N}$ iff every inner horn $\Lambda^n_k \to X$ ($0 < k < n$) extends \emph{uniquely} to $\Delta^n \to X$.}
\qaitem{Who is Grothendieck?}{Alexander Grothendieck (1928--2014), one of the most influential mathematicians of the 20th century. Reorganized algebraic geometry through category theory; introduced $\infty$-groupoids and the homotopy hypothesis in \emph{Pursuing Stacks} (1983).}
\qaitem{What changes when you drop ``uniquely''?}{You get quasi-categories (existence-only inner-horn filling). Composition becomes existential; associativity holds up to coherent homotopy. This is the foundation of $\infty$-category theory.}
\qaitem{Define a Kan complex.}{A simplicial set in which every horn $\Lambda^n_k \to X$ — inner or outer — extends to $\Delta^n \to X$.}
\qaitem{Who is Kan?}{Daniel M.\ Kan (1927--2013), Israeli-American mathematician at MIT. Introduced the horn-filling condition in 1958; also developed Kan extensions, adjoint functors, and the Kan model structure.}
\qaitem{What is an $\infty$-groupoid?}{A categorical structure in which all morphisms are invertible at every level (1-morphisms, 2-morphisms between them, and so on). Realized by Kan complexes.}
\qaitem{Where does ``groupoid'' come from?}{Heinrich Brandt (1927) coined the term for a generalization of groups in which composition is only partially defined.}
\qaitem{Is $\mathrm{Sing}(T)$ a Kan complex?}{Yes — because $\lvert\Lambda^n_k\rvert$ is a deformation retract of $\lvert\Delta^n\rvert$, every map from $\lvert\Lambda^n_k\rvert$ extends to $\lvert\Delta^n\rvert$, and the adjunction transports this to $\mathrm{Sing}(T)$.}
\qaitem{Is the nerve of an ordinary category always a Kan complex?}{No — outer horns fail in general. $\Lambda^2_0 \to \mathrm{N}(\mathcal{C})$ is two arrows out of a common source; filling forces an arrow relating their targets, which exists only if morphisms are invertible.}
\qaitem{When IS $\mathrm{N}(\mathcal{C})$ a Kan complex?}{When $\mathcal{C}$ is a groupoid (every morphism invertible). The intersection ``nerve $\cap$ Kan complex'' is exactly the image of $\mathbf{Grpd} \hookrightarrow \mathbf{sSet}$.}
\qaitem{How do nerve, quasi-category, and Kan complex relate?}{Nerve $\subset$ quasi-category (drop ``unique''); Kan complex $\subset$ quasi-category (specialize to all horns). Nerve of a groupoid = intersection of the two.}
\qaitem{What is the ``homotopy hypothesis''?}{Grothendieck's thesis ($\sim$ 1983) that $\infty$-groupoids are equivalent to homotopy types of spaces — i.e., the homotopy theory of Kan complexes is equivalent to the homotopy theory of topological spaces, by the realization–singular adjunction.}
\qaitem{Why is composition in a quasi-category not strict?}{Because the inner-horn filler gives \emph{a} composite, not \emph{the} composite — different choices of $\Delta^2$ filling a $\Lambda^2_1$ give different witnesses, all linked by higher-dimensional simplices.}
\qaitem{What category is $\tau_1 \mathrm{Sing}(T)$?}{The fundamental groupoid $\Pi_1(T)$: objects are points, morphisms are homotopy classes of paths, all invertible.}
\qaitem{Why do face operators on a nerve include composition?}{Because the $i$-th face for $0 < i < n$ drops the object at position $i$ from the chain, and the only way to restore composability is to compose the two adjacent morphisms — exactly what $f_{i+1} \circ f_i$ does.}
\qaitem{Is the realization functor a left adjoint? Why?}{Yes, it preserves colimits by construction (a coend is a colimit). The right adjoint is $\mathrm{Sing}$.}
\qaitem{What is the unit $X \to \mathrm{Sing}\,\lvert X\rvert$?}{The map sending an $n$-simplex $x \in X_n$ to the topological $n$-simplex it generates in $\lvert X\rvert$, regarded as a continuous map $\lvert\Delta^n\rvert \to \lvert X\rvert$.}
\end{qa}
\end{exercisesbox}


% ═══════════════════════════════════════════════════════════════════════════
\section*{Chapter Summary}
\addcontentsline{toc}{section}{Chapter Summary}
% ═══════════════════════════════════════════════════════════════════════════

\begin{summarybox}
\textbf{Geometric realization.}\quad $\lvert X\rvert = \int^{[n]} X_n \cdot
\lvert\Delta^n\rvert$, the coend that glues topological $n$-simplices
along simplicial face/degeneracy data. Left adjoint to $\mathrm{Sing}$;
left Kan extension of $\lvert\Delta^\bullet\rvert : \Delta \to \mathbf{Top}$
along Yoneda.

\textbf{The nerve.}\quad $\mathrm{N}(\mathcal{C})_n = \mathbf{Cat}([n], \mathcal{C})$,
the simplicial set of composable chains in $\mathcal{C}$. Fully faithful
$\mathbf{Cat} \hookrightarrow \mathbf{sSet}$, with left adjoint $\tau_1$
(fundamental category; ``$\tau$'' from \emph{truncation}).

\textbf{Grothendieck's characterization.}\quad A simplicial set is a
nerve iff every inner horn fills \emph{uniquely}. Relaxing ``uniquely''
to ``exists'' gives quasi-categories.

\textbf{Kan complexes.}\quad Simplicial sets in which every horn —
inner or outer — fills. Named for Daniel Kan (1958). $\mathrm{Sing}(T)$
is always a Kan complex. Kan complexes are $\infty$-groupoids:
outer-horn filling forces invertibility.

\textbf{The homotopy hypothesis.}\quad Grothendieck's thesis (1983):
$\infty$-groupoids and homotopy types of spaces are equivalent. Realized
by the realization–singular Quillen equivalence (Chapter~42).
\end{summarybox}


% ═══════════════════════════════════════════════════════════════════════════
\section*{Formal Restatement}
\addcontentsline{toc}{section}{Formal Restatement}
% ═══════════════════════════════════════════════════════════════════════════

\begin{formalrestatement}
\textbf{Realization.}\quad $\lvert{-}\rvert : \mathbf{sSet} \to \mathbf{Top}$,
$\lvert X\rvert = \int^{[n]} X_n \cdot \lvert\Delta^n\rvert
= \operatorname{Lan}_y \lvert\Delta^\bullet\rvert$.

\medskip
\textbf{Adjunction.}\quad $\lvert{-}\rvert \dashv \mathrm{Sing}$.
$\mathrm{Sing}(T)_n = \mathbf{Top}(\lvert\Delta^n\rvert, T)$.

\medskip
\textbf{Nerve.}\quad $\mathrm{N} : \mathbf{Cat} \to \mathbf{sSet}$,
$\mathrm{N}(\mathcal{C})_n = \mathbf{Cat}([n], \mathcal{C})$, with left
adjoint $\tau_1 = \operatorname{Lan}_y [-]$. $\mathrm{N}$ is fully faithful.

\medskip
\textbf{Characterization of the nerve's image.}\quad
\begin{equation*}
  X \in \mathrm{ess.im}(\mathrm{N}) \iff \forall n \geq 2, \, 0 < k < n, \,
  \forall \alpha : \Lambda^n_k \to X, \, \exists ! \, \tilde\alpha : \Delta^n \to X
  \text{ extending } \alpha.
\end{equation*}

\medskip
\textbf{Three horn conditions.}\quad For $X \in \mathbf{sSet}$:
\begin{itemize}
  \item \emph{Nerve}: inner horns fill uniquely.
  \item \emph{Quasi-category}: inner horns fill (Chapter~43).
  \item \emph{Kan complex}: all horns ($0 \leq k \leq n$, $n \geq 1$) fill.
\end{itemize}

\medskip
\textbf{Theorem.}\quad $\mathrm{Sing}(T)$ is a Kan complex for any $T \in
\mathbf{Top}$.

\medskip
\textbf{Inclusions.}\quad
$\{\mathrm{N}(\mathbf{Grpd})\} \;=\;
\{\mathrm{N}(\mathbf{Cat})\} \,\cap\, \{\mathrm{Kan}\} \;\subset\;
\{\mathrm{Kan}\} \;\subset\; \{\mathrm{quasi}\text{-}\mathrm{cat}\} \;\subset\; \mathbf{sSet}$.

\medskip
\textbf{The homotopy hypothesis (preview).}\quad $\mathrm{Kan} \simeq_{\mathrm{Q}}
\mathbf{Top}$: the homotopy theory of Kan complexes coincides with the
homotopy theory of topological spaces. Made precise by the Quillen
equivalence of Chapter~42.
\end{formalrestatement}
